% !TeX spellcheck = en_US
\documentclass[12pt]{article}

\input{common/variables}
\input{common/dates}

\setcounter{secnumdepth}{4}

\usepackage{times,fullpage,xspace,fancyhdr,url,color}
\usepackage[pdftex]{graphicx}
\usepackage[pdftex,
colorlinks=true,
urlcolor=black,
linkcolor=black,
citecolor=black,
bookmarksopen=false,
bookmarksnumbered=true,
pdfstartview=FitH]{hyperref}

\pdfcompresslevel=9

\hypersetup{
	pdftitle={\leaguename (\leaguenameabbr) Competition Rules},
	pdfauthor={Technical Committee \leaguenameabbr},
}
\usepackage{microtype}
\usepackage[utf8]{inputenc}
\usepackage{amsmath}
\usepackage{xargs}
\usepackage[colorinlistoftodos,prependcaption,textsize=tiny]{todonotes}
\usepackage{siunitx}
\usepackage[capitalize,noabbrev]{cleveref}
\usepackage[official]{eurosym}
\usepackage[useregional]{datetime2}
\usepackage{subcaption}
\usepackage{enumitem}
\usepackage{xcolor}
\DTMlangsetup[en-GB]{ord=raise,monthyearsep={,\space}}

% modern and easy to use tables tblr
\usepackage{tabularray}

\newcommandx{\unsure}[2][1=]{\todo[linecolor=red,backgroundcolor=red!25,bordercolor=red,#1]{#2}}
\newcommandx{\change}[2][1=]{\todo[linecolor=blue,backgroundcolor=blue!25,bordercolor=blue,#1]{#2}}
\newcommandx{\info}[2][1=]{\todo[linecolor=green,backgroundcolor=green!25,bordercolor=green,#1]{#2}}
\newcommandx{\improvement}[2][1=]{\todo[linecolor=orange,backgroundcolor=orange!25,bordercolor=orange,#1]{#2}}

% comment 'disable' in to disable all the todo notes :)
\usepackage
[
%disable
]{todonotes}

\usepackage{tikz-3dplot}

\usepackage{tocloft}
\cftsetindents{section}{0em}{2em}
\cftsetindents{subsection}{0em}{2em}
\renewcommand\cfttoctitlefont{\hfill\Large\bfseries}
\renewcommand\cftaftertoctitle{\hfill\mbox{}}

\sloppy
\newcommand{\ie}{\mbox{i.\,e.}\xspace}
\newcommand{\eg}{\mbox{e.\,g.}\xspace}
\newcommand{\cf}{see\xspace}
\newcommand{\inparagraph}[1]{\paragraph{#1\hspace{-1em} }}


% some colors
\definecolor{orange}{rgb}{1,0.5,0}
\definecolor{red}{rgb}{1,0,0}
\definecolor{green}{rgb}{0,1,0}


%%%%%%%%%%%%%%%%%%% GENERAL SECTION NUMBERING %%%%%%%%%%%%%%%%%%%%%%%

% Section numbering formats
\renewcommand{\thesection}{\Roman{section}}
\renewcommand{\thesubsection}{\arabic{subsection}}
\renewcommand{\thesubsubsection}{\arabic{subsection}.\arabic{subsubsection}}

% Counter hierarchy (important!)
\makeatletter
\@addtoreset{subsection}{section}
\@addtoreset{subsubsection}{subsection}
\makeatother

% unnumbered subsection that is still listed in toc
\makeatletter
\newcommand{\subsectionNoNumber}[1]{%
	\subsection*{#1}%
	\phantomsection%
	\addcontentsline{toc}{subsection}{\protect\numberline{}#1}%
}
\makeatother

%%%%%%%%%%%%%%%%%%% TEXT FORMAT %%%%%%%%%%%%%%%%%%%%%%%

\setlength{\parindent}{0pt}
\setlength{\parskip}{12pt plus 6pt minus 3 pt}
\widowpenalty=10000
\clubpenalty=10000

% needed to align an image and text correctly side by side
%\newcommand{\imagebox}[1]{\raisebox{2ex}{\raisebox{-\height}{#1}}}


%%%%%%%%%%%%%%%%%%% TITLE %%%%%%%%%%%%%%%%%%%%%%%

\title{Competition Rules\\\leaguename (\leaguenameabbr)}

\author{\leaguename Technical Committee\\\url{https://hsl.robocup.org/}}
\date{(DRAFT \RCYear Working Rules Document, as of \today)}

\renewcommand{\headrulewidth}{0.4pt}
\renewcommand{\footrulewidth}{0.4pt}

%%%%%%%%%%%%%%%%%%% PAGE HEADER FORMAT %%%%%%%%%%%%%%%%%%%%%%%

\pagestyle{fancy}
% make sure that a separate page can be identified 
\lhead{Competition Rules}
\chead{\leaguename}
\rhead{\today}
\lfoot{}
%\cfoot{}
%\cfoot{\thepage}
\rfoot{}


\begin{document}
	
	\pagenumbering{Roman}
	
	\maketitle
	\thispagestyle{empty}
	
	\begin{center}
		Questions or comments on these rules should be submitted via the following channels:
		\begin{itemize}
			\item \url{https://github.com/RoboCup-HumanoidSoccerLeague/HSL-Rules/issues};
			\item Discord server: \url{https://discord.gg/67upw6Nn}, channel \texttt{\#rule-book};
			\item Email: \url{hsl_tc@lists.robocup.org}
		\end{itemize}
	\end{center}
	
	\newpage
	\setcounter{tocdepth}{2}
	\tableofcontents
	
	\thispagestyle{fancy}
	
	\clearpage
	\pagenumbering{arabic}
	\setcounter{page}{1}
	
	\newpage
	
	\section{Purpose and Scope of this Document}

	This document details all rules and regulations required to run a RoboCup competition beyond gameplay related rules. This includes topics such as robot inspections, tropies and certificates, conflict resolution and more.
	
	\section{Robot Inspection}
	
	Every team participating in a RoboCup HSL competition must bring all robot platforms they intend to use to the league organizing committees for inspection before they can be used in a match. If a team has multiple different robot platforms or configurations, the team must bring one robot of each configuration to the inspection. If the team has multiple of the same robot platform, it is sufficient to only inspect one of them. The team must also bring one sample of each jersey they intend to use to the inspection.
	
	\subsection{Procedure}
	During inspection, the robot's height and weight are measured and compared against the allowed measurements as detailed in §3.5 of the rule book. Measurements are taken with the robot standing upright with its knees fully extended. The head of the robot must be oriented in such a way that it is tilted to either its maximum upwards tilt angle or the horizon line, whichever is lower. 
	
	The inspector then verifies that used sensors and other special modifications are in line with the rule book. This step is intended to identify clearly non-compliant features, an in-depth technical review will not be conducted. Where compliance cannot be determined from visual inspection alone, inspectors may request clarification from the team. The robot jerseys are also inspected.
	
	Finally, robots are inspected for potential safety problems and the team demonstrates the emergency stop procedure. The team must demonstrate that they are able to immediately and physically bring the robot into a safe state.
	
	\subsection{Pre-Approved Robots}
	Robot platforms that are included in the list of pre-approved standard platforms do not need to be measured. However, they still need to be brought to the inspection to verify that they are unmodified and that the emergency stop procedures work.
	
	\subsection{Reinspection}
	If a robot fails the initial inspection or is modified after inspection, it needs to be reinspected. Reinspection can be requested at any time. A robot without a valid inspection is not allowed to participate in any match. If a robot is modified from a previously approved version, only the modified components need to be inspected.
	
	\section{Referee Duty}
	Each team must name at least one person who is familiar with the rules and who might be assigned for referee duties and for the technical inspection by the league organizing committees.
	
	\section{Trophies and Certificates}
	
	A trophy is awarded to the winner, second and third place of the soccer tournament in each of the individual divisions. In case of less than 10 teams participating in a division the team ranked third will be awarded a certificate instead	of a trophy. In case of less than 6 teams participating in a division the team ranked second will also be awarded a certificate instead of a trophy.
	
	A trophy is awarded to the team ranked first on the technical challenges and certificates
	are awarded to the teams second and third in the technical challenges. In case less than 10 teams participate in the technical challenges in a division, the team ranked first in the technical challenges will be awarded a certificate instead of a trophy. 
	
	\todo{Best Humanoid award, to be discussed next meeting.}
	
	A certificate is awarded to the team that performed best in the open research challenge.
	
	A certificate is also awarded to the best referee as voted by the team leaders.
	
	\section{Forfeits}
	A team that forfeits is disqualified from the competition. Forfeiting is defined as refusing to make a good faith effort to participate in a scheduled game. Even if a team is unable to provide any running robots, their robot handlers and team leader are still expected to attend the match.
	
	If a team forfeits, the tournament structure is updated as follows:
	
	\begin{itemize}
		\item If a team chooses to forfeit a match in non-elimination games the other team plays on an empty goal.
		\item  If a team chooses to forfeit in the first round of knock-out games, it will be replaced by the runner up team.
		\item If a team chooses to forfeit in any other round of knock-out games, it will be replaced by the team that lost to the forfeiting team
		\item If a team chooses to forfeit the final game after the game for 3rd and 4th place has already begun, it will be replaced by the 3rd place winner instead. No new 4th place will be selected.
		\item   If a team chooses to forfeit the final before the game for 3rd and 4th place, it will be replaced by the team	that lost to the forfeiting team in the previous semi-finals. The team that lost to the forfeiting team in the previous semi-final.
	\end{itemize}

	A team forfeiting the final match should announce its decision at least 30 minutes before the start of the game for 3rd and 4th place. The league organization committee may impose a one year disqualification of the team and its members in case of avoidable delayed announcements.
	
	\section{Conflict Resolution}
	It is the responsibility of the team leader to inspect the other team’s robots an hour in advance of a game. Any concern regarding the rule compliance of any of the robots must be brought to the attention of the referee. If the referee is unavailable, they have to be brought to the attention of the Technical Committee instead.	
	
	Every result of a game needs to be signed by both team leaders and all robot handlers and referees involved in the game. Doubts concerning a serious violation of any rule during a specific game must be brought up to a member of the Technical Committee and investigated before signing the result. By signing the result sheet, a team agrees that the result arose from a fair game. Concerns must be brought to the attention of the Technical Committee within half an hour of the completion of the game. 
	
	If a team brings up an official concern to the Technical Committee, a meeting of the Technical Committee must be called as soon as possible. If the team of a member of the Technical Committee is	directly involved in the game in question, the respective member is excluded from the meeting. At least three members	of the Technical Committee need to be part of the meeting and the decision process. If less than three members of the Technical Committee are available, members of the Organizing committee or, if necessary, Trustees or members of committees from other leagues have to be called into the meeting. Members of these meetings may request to inspect the hardware, robot model and software of any team involved in the issue. If serious violations of rules or recurrent unsportive behavior are detected, the committee may, among others, decide to invalidate the result of the game in question or take disciplinary actions against a team as defined in Law 12.4, depending on the severity of the rule violation. The decisions of the committee meeting need to be announced to the whole league. If teams are requested to make changes to their code for the next game, they need to receive a period of at least four hours to make the requested change. If their next game was scheduled earlier than this, the game needs to be postponed.
\end{document}

